% !TEX root = ../Thesis.tex

\chapter{Introduzione}\label{ch:introduzione}

Garantire la confidenzialità, l'autenticità e l'integrità delle comunicazioni digitali è una delle esigenze fondamentali dell'architettura informatica contemporanea. Gran parte di questa garanzia è oggi affidata a primitive crittografiche e a protocolli consolidati come \emph{Transport Layer Security} (\textbf{TLS}), il cui scopo è rendere il contenuto dei messaggi inaccessibile a chiunque non sia in possesso delle chiavi appropriate. La crittografia moderna, in altre parole, protegge \emph{cosa} viene detto, ma non occulta il fatto che una comunicazione stia avvenendo.

Questa seconda proprietà, ovvero l'opacità del canale, è diventata cruciale in un numero crescente di scenari. La diffusione di meccanismi di sorveglianza di massa e di infrastrutture di censura, come quelle documentate rispetto al \emph{Great Firewall} cinese \cite{chinawall-2015} o relative al blocco di servizi di messaggistica in diversi contesti nazionali \cite{conger-2016, sablah-2025}, ha reso evidente che il traffico cifrato è facilmente identificabile e bloccabile da sistemi di filtraggio che operano a livello di rete. L'elevata entropia dei flussi TLS è, paradossalmente, un marcatore che li rende \emph{visibili}. In scenari operativi ostili, quindi, la sola crittografia non è sufficiente: diventa necessario ricorrere a metodologie che rendano indistinguibile una comunicazione cifrata da un normale flusso di dati innocui.

È in questo contesto che si colloca la \textbf{steganografia}. La disciplina steganografica si dedica all'occultamento di informazioni all'interno di oggetti di copertura (\emph{cover objects}) apparentemente ordinari, in modo che un osservatore esterno non possa nemmeno stabilire se una comunicazione nasconda informazioni al suo interno. A differenza della crittografia, la cui sicurezza poggia sulla complessità computazionale dei problemi matematici sottostanti, la sicurezza steganografica si fonda sull'\textbf{indistinguibilità statistica} \cite{turch-2020}: un sistema steganografico è tanto più sicuro quanto più la distribuzione dei propri \emph{stego-object} -- gli artefatti contenenti il messaggio nascosto -- è simile a quella degli oggetti naturali prodotti dal medesimo processo comunicativo.

Le tecniche steganografiche classiche, come la manipolazione dei bit meno significativi (\emph{Least Significant Bit}, \textbf{LSB}) nelle immagini, si sono rivelate progressivamente vulnerabili all'evoluzione delle tecniche di steganalisi, che combinano analisi statistica e apprendimento automatico per rilevare perturbazioni anche minime \cite{apau-2024, subramanian-2021}. La ricerca recente ha quindi intrapreso un cambio di paradigma: anziché modificare un contenuto esistente per nascondervi un messaggio, si sfruttano \textbf{modelli generativi} basati su reti neurali profonde per \emph{progettare} ex novo il contenuto stesso in funzione del messaggio da occultare. Questo approccio, indicato in letteratura come \emph{Steganography without Embedding} (\textbf{SwE}) \cite{liu-2020, dong-2025}, sposta la steganografia dal piano della perturbazione impercettibile a quello della generazione indistinguibile.

In questo filone si inserisce \textbf{Meteor} \cite{kaptchuk-2021}, un algoritmo di steganografia a chiave simmetrica pensato per il dominio testuale. Meteor sfrutta un \emph{Large Language Model} (\textbf{LLM}) per creare un canale di comunicazione occulto: anziché modificare un testo esistente, \emph{guida} il processo di generazione del modello, pilotando la scelta dei token successivi in base ai bit del messaggio segreto. La sicurezza del sistema poggia sulla garanzia formale che lo stego-text così prodotto sia statisticamente indistinguibile dall'output naturale del modello. Meteor costituisce, a oggi, una delle realizzazioni più convincenti del paradigma SwE.

La presente tesi si propone di investigare se un approccio concettualmente analogo a Meteor possa essere applicato al dominio visivo delle immagini. L'obiettivo non è implementare un sistema pronto all'uso, ma esplorare sistematicamente le condizioni sotto le quali un modello generativo visuale può essere impiegato come sampler steganografico, identificando i fattori strutturali che lo rendono possibile o, come si vedrà, lo limitano. L'indagine ruota attorno a due aspetti complementari: da un lato la scelta del tokenizer visivo, che determina la granularità della rappresentazione e la fedeltà del round-trip tra \emph{spazio latente} dei token e spazio dei pixel; dall'altro la \emph{natura autoregressiva} dei moderni generatori visivi, che introduce una dipendenza sequenziale tra token consecutivi e condiziona il modo in cui gli errori di trasmissione possono propagarsi.

Il percorso sperimentale, descritto in dettaglio nel \autoref{ch:sperimentazione}, si articola in tappe successive. Si parte da \textbf{VQ-GAN} \cite{esser-2020}, un'architettura rappresentativa della prima generazione di tokenizer visivi, per poi spostarsi su \textbf{Open-MAGVIT2} \cite{luo-2024}, che introduce una forma di quantizzazione (\emph{Lookup-Free Quantization}, \textbf{LFQ}) con proprietà strutturalmente diverse. Per ciascuna architettura, la sperimentazione valuta la fattibilità di guidare il campionamento autoregressivo come accade in Meteor, ricercando le cause che impediscono, allo stato attuale, l'applicazione di tale approccio al dominio visivo.

\bigskip

Il presente elaborato è organizzato come segue:
\begin{itemize}
	\item Il \textbf{\autoref{ch:fondamenti}} introduce i concetti teorici necessari: il modello di comunicazione steganografica, la nozione di indistinguibilità statistica, il paradigma SwE, il framework Meteor e le architetture di tokenizzazione visiva (VQ-GAN e Open-MAGVIT2 con LFQ) impiegate nella sperimentazione.

	\item Il \textbf{\autoref{ch:sperimentazione}} costituisce il nucleo del lavoro: descrive il percorso sperimentale seguito, i problemi incontrati, le soluzioni parziali esplorate e i risultati ottenuti. In particolare, viene analizzata la non-invertibilità della VQ-GAN, l'equivalenza concettuale dell'algoritmo di Meteor a livello di Transformer, il problema della propagazione dell'errore nei modelli autoregressivi e la strategia di embedding nei piani di bit.

	\item Il \textbf{\autoref{ch:conclusioni}} sintetizza i contributi della tesi, discute le limitazioni emerse e individua possibili direzioni di sviluppo futuro.
\end{itemize}